\section{Two-Phase Pretraining：高质量数据应该何时出现}

第二项数据研究把问题从 source selection 推进到 schedule：既然高质量 token 稀缺，是否应在训练末期提高其比例？两阶段方案不是简单“先差后好”，而是先用 broad mixture 建立一般能力，再用经过验证的高质量 blend 改变后期梯度，同时避免过度重复。


\lecturefigure{slide-046.jpg}{Two-Phase Pretraining：用阶段化数据提高 LLM accuracy}{CS25 V5 official deck, p.~46}


\lecturefigure{slide-047.jpg}{优化 LLM pretraining data selection：质量、混合与 schedule}{CS25 V5 official deck, p.~47}


\lecturefigure{slide-048.jpg}{关键贡献：定义、评估、fine-grained 分析与可扩展性}{CS25 V5 official deck, p.~48}

\teachervoice{00:21:03--00:22:30，Steven 把工作描述为系统化验证 two-phase pretraining，并强调先在更小 token budget 上 prototype blends。实践提示：数据配方和模型超参一样需要 cheap proxy experiments。}

\subsection{形式化：phase 与 blend 是两个独立旋钮}

前一节给出 two-phase 的直觉，本节把它拆成可独立控制的变量。只有形式化后，研究者才能分清收益来自第二阶段数据本身、出现时间、训练长度还是额外 token。这个区分还决定实验预算：固定总 token 可以研究 schedule，固定 phase duration 可以研究 blend，而同时改变两者只会得到一个不可归因的 recipe。

令第一阶段数据分布为 $q_1$、步数为 $T_1$，第二阶段为 $q_2$、步数为 $T_2$，总目标可写为
\[
\mathcal{L}=\sum_{t=1}^{T_1}\mathbb{E}_{x\sim q_1}[\ell_t(x)]
+\sum_{t=T_1+1}^{T_1+T_2}\mathbb{E}_{x\sim q_2}[\ell_t(x)].
\]
$q_2$ 的质量、多样性和重复率，与 $T_2/(T_1+T_2)$ 的比例是不同变量；一次只改变一个，才能判断收益来自 blend 还是 duration。


\lecturefigure{slide-049.jpg}{Two-phase approach：广覆盖 phase 1 与高质量 phase 2}{CS25 V5 official deck, p.~49}


\lecturefigure{slide-050.jpg}{Data sources：Common Crawl、Wikipedia、书籍、论文与任务数据}{CS25 V5 official deck, p.~50}


\lecturefigure{slide-051.jpg}{Evaluation suite：推理、知识、数学、编程与常识}{CS25 V5 official deck, p.~51}


\lecturefigure{slide-052.jpg}{Model/training setup：固定架构与统一训练设置}{CS25 V5 official deck, p.~52}

\begin{knowledgebox}{Upsampling 与 effective repetition}
若某数据源原有 $N$ tokens，在训练中被采样 $rN$ 次，则 upsampling factor 为 $r$。$r>1$ 会提高该域梯度权重，却也增加 memorization 和 diminishing return 风险。数据量、采样概率与训练步数必须一起报告。
\end{knowledgebox}

\subsection{Blend strategy：高质量不等于单一来源}

形式化给出了 $q_2$，现在要问它由哪些 sources 构成。所谓 high-quality 并不是一个来源名，而是过滤、文档结构、事实密度、领域匹配和重复率的组合；一本高质量书与一段高质量代码对不同 benchmark 的贡献也不同。因此读 blend 表不能只看 quality label，要追踪类别占比、每类 token 的实际重复次数，以及改善是否跨知识、推理、数学和 code 同时出现。

第二阶段把 high、medium、low-quality web 与 curated sources 重新混合。读 blend 表时先看每类数据占比，再看任务指标是否同步改善：某 blend 可能提高知识却伤害 code，或提高平均分却缩小长尾覆盖。


\lecturefigure{slide-053.jpg}{Two-phase data blending：phase 1 广覆盖、phase 2 偏高质量}{CS25 V5 official deck, p.~53}


\lecturefigure{slide-054.jpg}{Phase 2 blending 的性能影响：不同配方的多任务结果}{CS25 V5 official deck, p.~54}

\begin{knowledgebox}{深读：Phase-2 表格怎样避免 average-score illusion}
多任务表的平均分会掩盖任务间冲突。先把指标按 knowledge、reasoning、math、code 与 commonsense 分组，观察某 blend 是整体平移，还是只在某一组获益；再比较 absolute gain 与 baseline variance，避免把微小波动写成稳定提升。若配方含更多 task-like curated data，还要检查 benchmark format contamination。饼图表示 sampling share，不是 unique data share：一个占比小但反复采样的数据源可能贡献大量梯度。最后把 gains 除以额外 compute 和 data curation cost，得到每单位成本收益。只有当多个任务组、多个 seed 和更大规模都保持优势时，才有理由说 schedule 改善了 general pretraining，而不是对评测做了局部适配。
\end{knowledgebox}

\teachervoice{00:22:30--00:23:11，讲者说所有测试的 phase-two blends 都优于 simple continuation，并在扩大 token/model 后保持优势。讲义保留“tested blends”这个范围，不外推成任意二阶段配方都有效。}

\subsection{过采样与阶段长度：更多并不单调更好}

Blend 确定之后，训练仍有一个时间轴问题：相同高质量 token 应该出现几次、集中在多长的尾部阶段。若只观察 final loss，很容易把 memorization 当成有效学习；更可靠的读法是同时看 held-out task、domain diversity 与 phase length curve。图中的 turning point 不是固定常数，而是提醒团队必须搜索 duration，而不是默认“越久越好”。

如果第二阶段把少量高质量数据反复看到很多遍，训练 loss 仍可能下降，但泛化收益会饱和。可用一个简化权衡表示：
\[
G(r)=B(r)-\lambda M(r),
\]
其中 $B(r)$ 是更高质量曝光带来的收益，$M(r)$ 是重复导致的记忆/分布收窄成本，$\lambda$ 表示任务对过拟合的敏感度。最优 $r$ 通常位于内部而非无限大。


\lecturefigure{slide-055.jpg}{Phase 2 upsampling 是否有效：收益、占比与递减回报}{CS25 V5 official deck, p.~55}


\lecturefigure{slide-056.jpg}{Two-phase scalability：更大 token budget 与模型下的趋势}{CS25 V5 official deck, p.~56}


\lecturefigure{slide-057.jpg}{Phase duration：后期高质量阶段太短或太长都可能非最优}{CS25 V5 official deck, p.~57}

\begin{knowledgebox}{深读：Phase duration 为什么会改变模型而不仅是数据权重}
训练早期模型表示还在形成，广覆盖数据能建立词汇、事实和通用模式；训练后期 learning rate、optimizer state 与参数所在 basin 已不同，同一批高质量 token 在此时产生的更新方向也不同。因此 two-phase schedule 不是把两个静态数据集简单加权，它还利用了优化路径依赖。读 duration 曲线时应同时报告 learning-rate schedule 和是否重置 optimizer；若 phase 2 恰逢学习率变化，收益可能由两者共同造成。还要检查被 phase 2 提升的任务是否牺牲 multilingual、creative 或 long-tail 能力。工程上可以把 phase duration 做成 Pareto search：横轴为高质量 unique tokens/repetitions，纵轴为多任务收益和遗忘，再选择稳定区域，而不是峰值单点。
\end{knowledgebox}

\teachervoice{00:23:11--00:24:00，课堂强调 phase 2 不能无限拉长，过多 upsampling 会产生 detrimental 或 diminishing returns。这个 caveat 比“高质量数据有效”更接近可执行配方。}

\begin{lstlisting}[caption={Two-phase data schedule 的可复现实验骨架}]
phase1 = mix(broad_web, books, code, multilingual)
phase2_candidates = propose_quality_blends()
for blend in phase2_candidates:
    proxy = train_small_model(phase1, blend, fixed_tokens=True)
    evaluate(proxy, multi_task_suite, multiple_seeds=True)
selected = choose_pareto_blend(quality, diversity, stability)
train_full_model(phase1, selected, report_repetition=True)
\end{lstlisting}

\subsection{把两项研究合成 data strategy}

TinyDialogues 改变 source structure；two-phase pretraining 改变 schedule。二者共同说明，data effectiveness 不是样本的单一内在分数，而是 source、model、objective、训练阶段和评测之间的匹配关系。


\lecturefigure{slide-058.jpg}{Two-phase findings：结构化高质量数据与可扩展性的综合结论}{CS25 V5 official deck, p.~58}


\lecturefigure{slide-060.jpg}{Unified insights：小数据与大规模 pretraining 的共同 data lessons}{CS25 V5 official deck, p.~60}


\lecturefigure{slide-061.jpg}{数据研究总结：质量、结构、调度与可复现实验同等重要}{CS25 V5 official deck, p.~61}

\begin{knowledgebox}{术语消化：source、blend、schedule 与 curriculum}
\textbf{Source} 是原始数据类别；\textbf{blend} 是同一阶段的采样比例；\textbf{schedule} 描述比例随训练步数怎样变化；\textbf{curriculum} 还可能按难度或能力有意排序。Two-phase pretraining 是一种 piecewise-constant schedule，不等同于所有 curriculum。读 slides 53--61 时要同时报告 phase token 数、每类数据 unique tokens、effective repetitions 和模型规模，否则无法复现“高质量阶段”的真实强度。
\end{knowledgebox}

\begin{importantbox}{Data-centric decision rule}
先问“缺什么能力”，再选择能提供对应监督的 source；用小规模多 seed 实验筛 blend；扩大规模时检查排序是否保持；最后报告重复率、版权/隐私与 long-tail 代价。不要把一张 aggregate score 表当作 universal recipe。
\end{importantbox}

\subsection{本章小结}

Two-phase pretraining 的核心不是“最后喂好数据”，而是把 blend 与 duration 变成可消融变量。高质量阶段可以改善多任务性能，但过采样存在递减回报；最可靠的工程路径是 proxy experiment、multi-task evaluation 和 scale validation。

